\documentclass[11pt]{article}
\usepackage[margin=1in]{geometry}
\usepackage{amsmath,amssymb,amsthm}
\usepackage[hidelinks]{hyperref}

\newtheorem{theorem}{Theorem}
\newtheorem{lemma}[theorem]{Lemma}
\newcommand{\Tr}{\operatorname{Tr}}
\newcommand{\KF}{K}

\title{Compatible Spin Alignment for Arbitrary Qubit Reference States\\
\large A three-qubit, pair-support theorem}
\author{Draft research note --- authorship and venue pending}
\date{August 2026}

\begin{document}
\maketitle

\begin{abstract}
We prove the compatible-marginal spin-alignment majorization inequality for
three qubits, arbitrary qubit reference state, and arbitrary probability
weights supported on the three two-body subsets.  The proof has two parts.  We
first establish the maximally mixed endpoint from the published polygon
inequalities for pure qubit marginals, a two-projection calculation, and a
spin-flip identity.  We then interpolate to every qubit reference state using
Ky Fan convexity and a common ordered target spectrum. The endpoint argument
substantially overlaps with the earlier Song--Chen preprint; no novelty claim
is made for that endpoint.
\end{abstract}

\section{Context and conventions}

Song and Chen formulate the compatible-marginal spin-alignment conjecture for
a global state, its compatible marginals, an arbitrary subset measure, and a
reference state \(Q\) \cite[Conjecture 2]{SongChen}. Their Proposition 3
handles the three-qubit, pair-support \(Q=I/2\) endpoint, and their outlook
asks for arbitrary positive semidefinite \(Q\) in more general three-qubit
settings. This note proves the qubit, same-pair-support extension.

For a Hermitian matrix \(X\), write
\(\lambda_1(X)\ge\cdots\ge\lambda_8(X)\) for its decreasing eigenvalues and
\(\KF_r(X)=\sum_{i=1}^r\lambda_i(X)\). We write \(X\preceq Y\) when all Ky
Fan sums of \(X\) are at most the corresponding sums of \(Y\), with equal
trace. The distinct notation \(X\succeq Y\) means that \(X-Y\) is positive
semidefinite.

\section{Statement}

For a three-qubit state \(\rho_{ABC}\), let \(\rho_{AB},\rho_{AC},\rho_{BC}\)
be its marginals.  Let \(Q\) be a qubit state and choose a maximal eigenvector
\(|q_1\rangle\).  For \(a,b,c\ge0\), \(a+b+c=1\), define
\[
\begin{aligned}
H_Q={}&a\rho_{AB}\otimes Q_C+b\rho_{AC}\otimes Q_B+
cQ_A\otimes\rho_{BC},\\
T_Q={}&a|q_1q_1\rangle\!\langle q_1q_1|_{AB}\otimes Q_C+
b|q_1q_1\rangle\!\langle q_1q_1|_{AC}\otimes Q_B+
cQ_A\otimes|q_1q_1\rangle\!\langle q_1q_1|_{BC}.
\end{aligned}
\]
The \(AC\) terms are in ambient \(ABC\) order.

\begin{theorem}
For every \(\rho_{ABC}\), every probability vector \((a,b,c)\), and every
qubit state \(Q\), one has \(H_Q\preceq T_Q\).
\end{theorem}

This is the compatible-marginal conjecture only for three qubits and measure
supported on \(AB,AC,BC\).  It does not cover incompatible triples, other
supports, more parties, or higher local dimension.

\section{The \(Q=I/2\) endpoint}

\begin{lemma}\label{lem:endpoint}
For every three-qubit state \(\rho_{ABC}\) and \(a,b,c\ge0\) with
\(a+b+c=1\),
\[
a\rho_{AB}\otimes I_C+b\rho_{AC}\otimes I_B+cI_A\otimes\rho_{BC}
\preceq
aP_{00,AB}\otimes I_C+bP_{00,AC}\otimes I_B+cI_A\otimes P_{00,BC}.
\tag{1}
\]
\end{lemma}

\begin{proof}
Permute systems so \(a\ge b\ge c\).  Let \(X_1,X_2,X_3\) be the successive
one-, two-, and three-term sums on the left of (1), and let \(T_1,T_2,T_3\)
be their aligned versions.  Then the two sides of (1) are
\[
(a-b)X_1+(b-c)X_2+cX_3,\qquad
(a-b)T_1+(b-c)T_2+cT_3. \tag{2}
\]

The spectra of \(T_1,T_2,T_3\) are respectively
\[
(1,1,0^6),\qquad(2,1,1,0^5),\qquad(3,1,1,1,0^4). \tag{3}
\]
The first prefix is immediate: \(X_1=\rho_{AB}\otimes I_C\) has first two
Ky Fan sums at most \(1,2\).

For the second prefix, first take two pure pair states.  Their extensions to
\(ABC\) are rank-two projections \(P=VV^\dagger\) and \(R=WW^\dagger\).
If \(s_1\ge s_2\) are the singular values of \(V^\dagger W\), then their
four principal-angle eigenvalues in \(P+R\) are
\[
1+s_1,\quad1+s_2,\quad1-s_2,\quad1-s_1. \tag{4}
\]
Writing the normalized pair-state coefficient matrices as \(\Phi,\Psi\),
\[
|\phi\rangle=\sum_{a,b}\Phi_{ab}|a,b\rangle,\quad
|\psi\rangle=\sum_{a,c}\Psi_{ac}|a,c\rangle,\quad
V|c\rangle=|\phi\rangle_{AB}|c\rangle_C,\quad
W|b\rangle=|\psi\rangle_{AC}|b\rangle_B.
\]
Thus \((V^\dagger W)_{c,b}=(\Phi^\dagger\Psi)_{b,c}\), hence
\[
s_1+s_2=\|\Phi^\dagger\Psi\|_1\le\|\Phi\|_F\|\Psi\|_F=1. \tag{5}
\]
Thus the first three Ky Fan sums are at most \(2,3,4\), as in (3).
Independent spectral decompositions and Ky Fan convexity extend this result
to mixed pair states, proving \(X_2\preceq T_2\).

For the third prefix, Ky Fan convexity reduces each index to a pure global
state.  Independently diagonalize its one-qubit marginals as
\[
\rho_i=\operatorname{diag}(1-r_i,r_i),\qquad0\le r_i\le\tfrac12.
\]
Higuchi--Sudbery--Szulc prove the polygon inequalities
\(r_i\le r_j+r_k\) \cite[Eq.~(2), p.~1]{Higuchi}.  With
\[
s=r_A+r_B+r_C,\quad\delta=(1-s)_+,\quad
D=\rho_A\otimes I\otimes I+I\otimes\rho_B\otimes I+
I\otimes I\otimes\rho_C-I,
\]
the qubit identity
\(\sigma_yM^T\sigma_y=(\Tr M)I-M\) gives
\[
X_3=D+\rho+\widetilde\rho,\qquad
\widetilde\rho=(\sigma_y^{\otimes3})\rho^T(\sigma_y^{\otimes3})\succeq0.
\tag{6}
\]
The diagonal entries of \(D\) are
\[
2-s,\ 1-s+2r_C,\ 1-s+2r_B,\ 1-s+2r_A,\
s-2r_A,\ s-2r_B,\ s-2r_C,\ s-1. \tag{7}
\]
If \(h_1\ge\cdots\ge h_8\) denotes their decreasing rearrangement, the
polygon inequalities first show that every entry in (7) is at least \(s-1\):
for the three displayed entry types, the differences are bounded below by
\(3-2s\), \(2(1-r_j-r_k)\), and \(1-2r_i\), respectively. Hence
\(h_8=s-1\). If \(s\ge1\), every entry in (7) is at most one. If
\(s\le1\), only \(2-s=1+\delta\) can exceed one: the polygon inequalities
bound \(1-s+2r_i\le1\), and all remaining entries are at most one. Therefore
\[
h_8=s-1,\qquad\KF_r(D)\le r+\delta\quad(r=1,2,3). \tag{8}
\]
Since \(X_3\succeq D\) and \(X_3\succeq0\),
\[
\sum_{i=r+1}^{8}\lambda_i(X_3)
\ge4-\KF_r(D)+\delta,
\]
so \(\KF_r(X_3)\le2+\KF_r(D)-\delta\le r+2\) for \(r=1,2,3\).
Positivity and trace handle \(r\ge4\).  Therefore \(X_3\preceq T_3\).

Ky Fan subadditivity in (2) proves (1), because the aligned targets have a
common nested diagonal order and the right-hand side has spectrum
\((1,a,b,c,0^4)\).
\end{proof}

\section{Interpolation to arbitrary \(Q\)}

\begin{proof}[Proof of the theorem]
Apply one common local unitary sending the chosen maximal eigenvector to
\(|0\rangle\) and diagonalizing \(Q\) as
\(Q_q=\operatorname{diag}(q,1-q)\), \(1/2\le q\le1\).  This preserves both
compatibility and majorization, including the degenerate case \(q=1/2\).
Set \(t=2q-1\).  Since each selected support has one complement qubit,
\[
Q_q=(1-t)I/2+tP_0,\quad
H_q=(1-t)H_{1/2}+tH_1,\quad
T_q=(1-t)T_{1/2}+tT_1. \tag{9}
\]
Lemma~\ref{lem:endpoint}, scaled by \(1/2\), gives
\(H_{1/2}\preceq T_{1/2}\). At \(q=1\),
\(H_1\) is a density matrix and \(T_1=P_{000}\), hence
\(H_1\preceq T_1\).

After the earlier permutation of weights, the target spectrum is
\[
\lambda(T_q)=\bigl(q,(1-q)a,(1-q)b,(1-q)c,0^4\bigr). \tag{10}
\]
It remains in this order throughout the interval, so every target Ky Fan sum
is affine in \(t\).  Ky Fan convexity now gives
\[
\KF_k(H_q)\le(1-t)\KF_k(H_{1/2})+t\KF_k(H_1)
\le(1-t)\KF_k(T_{1/2})+t\KF_k(T_1)=\KF_k(T_q).
\]
Equal trace completes the proof.
\end{proof}

\section{Relation to prior work and replay}

Song and Chen state the same \(Q=I/2\) compatible endpoint in their current
preprint \cite[Proposition 3]{SongChen}; our endpoint derivation is an
independently certified reconstruction of substantially the same proof
architecture and is not claimed as new. The arbitrary-\(Q\) interpolation
answers the qubit, same-support portion of a direction named in their
outlook. The bounded literature/overlap assessment is recorded separately
with this paper phase.
Alhejji and Knill's classical-state result subsumes our earlier diagonal
consistency check \cite[Proposition IV.5]{AlhejjiKnill}.

The project-root replay materials named in REPLAY.md contain the immutable
C79 endpoint certificate, C78 reinstatement chain, and exact checks of the
spin-flip expansion, polygon polytope, target assembly, and interpolation
spectrum. A standalone release archive has not yet been built.

\begin{thebibliography}{9}
\bibitem{SongChen}
Z. Song and L. Chen,
\newblock \emph{A counterexample to the strong spin alignment conjecture},
\newblock arXiv:2603.25410 (2026).

\bibitem{Higuchi}
A. Higuchi, A. Sudbery, and J. Szulc,
\newblock \emph{One-qubit reduced states of a pure many-qubit state:
polygon inequalities},
\newblock Physical Review Letters 90, 107902 (2003).

\bibitem{AlhejjiKnill}
M. A. Alhejji and E. Knill,
\newblock \emph{Towards a resolution of the spin alignment problem},
\newblock Communications in Mathematical Physics 405, 119 (2024).
\end{thebibliography}

\end{document}
